\documentclass[10pt,a4paper]{amsart}
\usepackage[margin=19mm]{geometry}
\usepackage{amsmath,amssymb,mathtools}
\usepackage[scr=rsfs,cal=euler]{mathalfa}
\usepackage{xfrac}
\usepackage[unicode,hidelinks,bookmarks=false]{hyperref}
\newcommand{\ie}{i.e.\ }
\newcommand{\cf}{cf.\ }
\newcommand{\resp}{respectively\ }
\newcommand{\Subsection}{\subsection}
\providecommand{\nobreakdash}{\nobreak\hbox{-}}
\setlength{\emergencystretch}{2em}
\multlinegap0pt
\usepackage{bm}
\usepackage{bbm}
\usepackage{dsfont}
\usepackage{xspace}
\usepackage{cancel}
\usepackage{mathtools}
\usepackage{amsthm}
\makeatletter
\@ifundefined{assumption}{\newtheorem{assumption}{Assumption}}{}
\@ifundefined{lem}{\newtheorem{lem}{Lemma}}{}
\makeatother
\makeatletter
\renewcommand{\H}{\mathcal{H}}
\newcommand{\M}{\mathcal{M}}
\newcommand{\R}{\mathbb{R}}
\expandafter\ifx\csname bluetext\endcsname\relax
\newenvironment{bluetext}
  {\color{black}\ignorespaces}
\fi
\expandafter\ifx\csname proof\endcsname\relax
\renewenvironment{proof}%
  {\begin{trivlist}%
   \item[]\ignorespaces}%
\fi
\expandafter\ifx\csname redtext\endcsname\relax
\newenvironment{redtext}
  {\color{black}\ignorespaces}
\fi
\makeatother
\title[A Unified Framework for Wasserstein Convergence of ULMC Meth]{A Unified Framework for Wasserstein Convergence of ULMC Methods beyond Log-Concavity: Old and New}
\author{Wanjie Lyu, Xiaojie Wang, Bin Yang}
\date{Source: arXiv:2609.20713v1 (CC BY 4.0)}
\begin{document}
\maketitle

\noindent\emph{Source and licence.} A Unified Framework for Wasserstein Convergence of ULMC Methods beyond Log-Concavity: Old and New. Version arXiv:2609.20713v1 (CC BY 4.0), distributed under the Creative Commons Attribution 4.0 International licence, as recorded in the arXiv licence metadata for this exact version and shown on the abstract page https://arxiv.org/abs/2609.20713v1. Full rights notice: licenses/arxiv-2609.20713-CC-BY-4.0.txt.\par\smallskip

\noindent\textit{Editorial scope.} Counted unit: the Lem whose source label is \texttt{lem:velocity-one-step} in the article (source0.tex lines 7027--7146, source label \texttt{lem:velocity-one-step}), delivered together with the article's own complete original proof of it (source0.tex lines 7148--7681, one \texttt{proof} environment of 534 lines). No mathematics is written, repaired or re-derived: statement and proof are the article's own text line for line. Required context retained uncounted: as:Lip (lines 2534--2543); as:U\_0-bound (lines 2545--2557); eq:uni-process (lines 2743--2782); as:coeff-approx-general (lines 2960--2996); eq:bound-Psi-Phi (lines 3060--3065); lem:nabla-tau-delta (lines 4721--4779); eq:elementary\_inequality (lines 4829--4840); eq:gradient-growth-square (lines 4917--4924); eq:grad-Xtau-split-clear (lines 5216--5294). Every definition, display and label of those lines is retained unchanged. Omitted from this package and not redistributed: 1--2533, 2544--2544, 2558--2742, 2783--2959, 2997--3059, 3066--4720, 4780--4828, 4841--4916, 4925--5215, 5295--7026, 7147--7147, 7682--17225. Nothing inside a retained excerpt is omitted or altered: every retained body is delivered exactly as the article's lines are written, so no omission receipt of any kind accompanies this package. Citations: the retained excerpts carry no citation command at all, so no bibliography entry of the article is transcribed and no reference is redistributed. Reference adaptation: none was needed and none was made; every reference inside the delivered unit resolves inside it, so no unresolved reference or citation is produced. Printed slips kept literal: no printed word, symbol, hypothesis or number of the counted statement or of its proof was altered. Layout and class: the article's own class is \texttt{article} with packages including amssymb, multirow, makecell, graphicx, subcaption, caption, array, tabularray, amsmath, mathrsfs, amsfonts, indentfirst, colortbl, dcolumn; the wrapper is the lane's pinned \texttt{amsart} editorial wrapper at 10pt on A4, which restates the theorem-like environments the retained text needs and adds the article's own macro definitions that the retained text needs (\texttt{\textbackslash H}, \texttt{\textbackslash M}, \texttt{\textbackslash R}), with extra wrapper packages (bm, bbm, dsfont, xspace, cancel, mathtools). The delivered numbering is the unit's own. Assets: no retained span contains a figure environment, a graphics-inclusion command, a PDF-page inclusion, a table or a TikZ picture; no image file is copied into this package, the package declares no asset dependency and no image file is redistributed with it. Licence: version arXiv:2609.20713v1 of this article is distributed under the Creative Commons Attribution 4.0 International licence, as recorded in the arXiv licence metadata for this exact version and shown on the abstract page https://arxiv.org/abs/2609.20713v1. A full rights notice is delivered at licenses/arxiv-2609.20713-CC-BY-4.0.txt. Characters: every retained excerpt is pure ASCII, so the delivered engine, XeLaTeX with its default Latin Modern text font, reports no missing character.
\begin{assumption}[\textbf{Gradient Lipschitz}]\label{as:Lip}
The potential is continuously differentiable with an $L$-Lipschitz gradient, meaning that there exists a dimension-independent constant $L > 0$ such that
\begin{align}\label{eq:lip}
    |\nabla U(x) - \nabla U(y)| \leq L |x - y|,
    \quad 
    \forall
    x, y \in
    \mathbb{R}^d.
    \end{align}
\end{assumption}

\begin{assumption}\label{as:U_0-bound}
The potential $U$ is non-negative and satisfies 
\begin{align}
|U(0)|
\leq 
L_1, 
\quad 
|\nabla U(0)|
\leq 
L_2 d^{\frac{1}{2}},
\end{align}
where $L_1, L_2 > 0$ are constants independent of dimension $d$.
\end{assumption}

\begin{equation}
\label{eq:uni-process}
\left\{
\begin{aligned}
&\bar{X}_{k+\tau_k}
=
\bar{X}_k
+ 
\tau_k h\Phi_1(\tau_kh)\bar{V}_k
-
\ell_1 \alpha \gamma^{-1}
\tau_k h 
(1-\Phi_2(\tau_k h))
\nabla U (\bar{X}_k)
+
\ell_2\sqrt{2\gamma \alpha}
\Delta W_{k+\tau_k}^{\Phi_3,1}, \\
&\bar{X}_{k+1}
=
\bar{X}_k
+
h\Phi_1(h)\bar{V}_k
-
\alpha 
h(h-\tau_k h)\Phi_2(h-\tau_kh)\nabla U(\bar{X}_{k+\tau_k})
+
\sqrt{2\gamma \alpha}
\Delta W_{k+1}^{\Phi_3,1},\\
&\bar{V}_{k+1}
=
\Psi_1(h)\bar{V}_k
-
h\Psi_2(h-\tau_k h)\alpha
\nabla U(\bar{X}_{k+\tau_k})
+
\sqrt{2\gamma\alpha}
\Delta W_{k+1}^{\Psi_3,2},
\end{aligned}
\right.
\end{equation}

\begin{assumption}[\textbf{Consistency conditions on coefficients}]
\label{as:coeff-approx-general}
Let \(\boldsymbol{\eta}=(\eta_1,\dots,\eta_6)\) be a vector with each \(\eta_i\ge 0\).
% $$
% \boldsymbol{\eta}
% :=
% (\eta_1,\eta_2,\eta_3,\eta_4,\eta_5,\eta_6)
% \in [0,\infty)^6.
% $$
For any uniform stepsize
\(h\in(0,1\wedge\gamma^{-1}]\) and all \(r\in(0,h]\), assume that the coefficient functions in \eqref{eq:uni-process} satisfy
\begin{equation}
\label{eq:general-coeff-approx}
\begin{aligned}
|\Phi_1(r)-\phi(r)|
&\leq \kappa r^{\eta_1},
&
|\Phi_2(r)-\phi(r)|
&\leq \kappa r^{\eta_2},
&
|\Phi_3(r)-\phi(r)|
&\leq \kappa r^{\eta_3},
\\
|\Psi_1(r)-\psi(r)|
&\leq \kappa r^{\eta_4},
&
|\Psi_2(r)-\psi(r)|
&\leq \kappa r^{\eta_5},
&
|\Psi_3(r)-\psi(r)|
&\leq \kappa r^{\eta_6},
\end{aligned}
\end{equation}
where \(\kappa\geq0\) is independent of \(h\), \(r\), and the dimension $d$.
%%

\end{assumption}

\begin{align}\label{eq:bound-Psi-Phi}
    |\Psi_i(r)|\leq (1+\kappa),
    \quad |\Phi_i(r)| \leq (1+\kappa), \quad  i \in 
    \{1,2,3\},
    \quad \forall \, r\in(0,h\wedge\gamma^{-1}].
\end{align}

\begin{lem}\label{lem:nabla-tau-delta}
Let Assumptions \ref{as:Lip}, \ref{as:U_0-bound} hold and suppose that Assumption \ref{as:coeff-approx-general} is satisfied with $\boldsymbol{\eta}=\boldsymbol{\eta}_{\mathrm M}$. Let $\{(\bar{X}_k,\bar{V}_k)\}_{k\geq0}$ be generated by the U-ULMC scheme \eqref{eq:uni-process}. Then, for any uniform stepsize $h\in(0,1\wedge\gamma^{-1}]$, it holds that 
%%
%%
\begin{equation}
\begin{aligned}
  \mathbb{E}
\Big[\big|
\nabla U(\bar{X}_{k+\tau_k})
-
\nabla U(\bar{X}_{k})
\big|^2\Big]  
\leq   
\Big( 
\H^{x}_1
\mathbb{E}
\Big[\big|
\bar{X}_k
\big|^2\Big]
+
\H^{v}_1
\mathbb{E}
\Big[\big|
\bar{V}_k
\big|^2\Big]
+
\H_1
d
\Big)h^2,
\end{aligned}
\end{equation}
where 
\begin{equation}
\begin{aligned}
\H^{x}_1
:=   &
6L^4
\gamma^{-2}
\alpha^2
\ell_1^2
(2 +  \kappa)^2,
& &
\H^{v}_1
:=   
3L^2(1+\kappa)^2,
\\
\H_1
:=   &
6L^2L_2^2
\gamma^{-2}
\alpha^2
\ell_1^2
(2 +  \kappa)^2
+
2L^2
\gamma\alpha\ell_2^2(1+\kappa)^2.
\end{aligned}
\end{equation}
\end{lem}

\begin{equation}
\label{eq:elementary_inequality}
\Big(
\sum_{i=1}^{k}|u_i|
\Big)^2 
\leq 
k
\sum_{i=1}^{k}|u_i|^2, 
\quad
u_i \in \R, \; 
\forall k \in \mathbb{N},
\end{equation}

\begin{equation}\label{eq:gradient-growth-square}
|\nabla U(\bar{X}_k)|^2
\leq
2L^2
|\bar{X}_k|^2
+
2L_2^2d.
\end{equation}

\begin{equation}
\label{eq:grad-Xtau-split-clear}
\begin{aligned}
\mathbb{E}
\Big[
\big|
\nabla U(\bar{X}_{k+\tau_k})
\big|^2
\Big]
\leq   &
2
\mathbb{E}
\Big[
\big|
\nabla U(\bar{X}_k)
\big|^2
\Big]
+
2\mathbb{E}
\Big[
\big|
\nabla U(\bar{X}_{k+\tau_k})
-
\nabla U(\bar{X}_k)
\big|^2
\Big]
\\   
\leq   &
4 L^2
\mathbb{E}
\Big[
\big|
\bar{X}_k
\big|^2
\Big]
+
2
\Big( 
\H^{x}_1
\mathbb{E}
\Big[\big|
\bar{X}_k
\big|^2\Big]
+
\H^{v}_1
\mathbb{E}
\Big[\big|
\bar{V}_k
\big|^2\Big]
+
\H_1
d
\Big)h^2
+
4 L_2^2  d
\\ 
\leq   &
(4 L^2
+
2\H^{x}_1)
\mathbb{E}
\Big[\big|
\bar{X}_k
\big|^2\Big]
+
2 
\H^{v}_1
\mathbb{E}
\Big[\big|
\bar{V}_k
\big|^2\Big]
+
( 
2\H_1  +  
4 L_2^2 
)
d,
\end{aligned}
\end{equation}

\begin{lem}
\label{lem:velocity-one-step}
Let Assumptions \ref{as:Lip}, \ref{as:U_0-bound} hold and suppose that Assumption \ref{as:coeff-approx-general} is satisfied with $\boldsymbol{\eta}=\boldsymbol{\eta}_{\mathrm M}$. Let $\{(\bar{X}_k,\bar{V}_k)\}_{k\geq0}$ be generated by the U-ULMC scheme \eqref{eq:uni-process}. Then, for any uniform stepsize $h\in(0,1\wedge L^{-1}\wedge\gamma^{-1}]$, it holds that
\begin{equation}
\label{eq:velocity-one-step}
\begin{aligned}
\tfrac{1}{4}
\mathbb{E}
\Big[
\big|
\bar{V}_{k+1}
\big|^2
\Big]
\leq&
\tfrac{1}{4}
\mathbb{E}
\Big[
\big|
\bar{V}_{k}
\big|^2
\Big]
-
\tfrac{\gamma}{2}h
\mathbb{E}
\Big[
\big|
\bar{V}_{k}
\big|^2
\Big]
-
\tfrac{\alpha}{2}h
\mathbb{E}
\Big[
\big\langle
\bar{V}_k,
\nabla U(\bar{X}_k)
\big\rangle
\Big]
\\
&+
\Big(
\M^x_3
\mathbb{E}
\Big[
\big|
\bar{X}_k
\big|^2
\Big]
+
\M^v_3
\mathbb{E}
\Big[
\big|
\bar{V}_k
\big|^2
\Big]
\Big)h^2
+
\M_3dh,
\end{aligned}
\end{equation}
where
\begin{equation}
\label{eq:M3-constants}
\begin{aligned}
\M^x_3
:=&
\tfrac{1}{2}
(3\alpha^2+\alpha\gamma)L^2
+
\tfrac{9}{2}
\alpha^2(\gamma+\kappa)^2
\big(
2L^2+\H_1^x
\big)
+
\tfrac{9}{4}
\alpha^2\H_1^x,
\\
\M^v_3
:=&
\tfrac{1}{4}
(\gamma^2+\alpha\gamma)
+
\tfrac{1}{2}
+
\tfrac{9}{4}
\big(
\tfrac{\gamma^2}{2}+\kappa
\big)^2
+
\tfrac{9}{2}
\alpha^2(\gamma+\kappa)^2
\H_1^v
+
\tfrac{9}{4}
\alpha^2\H_1^v,
\\
\M_3
:=&
\tfrac{1}{2}
(3\alpha^2+\alpha\gamma)
L_2^2
+
\tfrac{9}{2}
\alpha^2(\gamma+\kappa)^2
\big(
2L_2^2+\H_1
\big)
+
\tfrac{9}{4}
\alpha^2\H_1
+
\gamma\alpha(1+\kappa)^2.
\end{aligned}
\end{equation}
%%
%%
%%
\end{lem}

\begin{proof}
\textbf{Proof.}
Recalling the velocity update in \eqref{eq:uni-process}, we have
\begin{equation}
\label{eq:velocity-update}
\bar{V}_{k+1}
=
\Psi_1(h)\bar{V}_k
-
\alpha h\Psi_2(h-\tau_kh)
\nabla U(\bar{X}_{k+\tau_k})
+
\sqrt{2\gamma\alpha}
\Delta W_{k+1}^{\Psi_3,2}.
\end{equation}
We first rewrite this update by separating the leading part from the coefficient errors. Since
\(
\psi(h)=e^{-\gamma h},
\)
we have, for $h\leq\gamma^{-1}$,
\[
|\psi(h)-(1-\gamma h)|
\leq
\tfrac{\gamma^2}{2}h^2.
\]
Together with Assumption \ref{as:coeff-approx-general}, this gives
\begin{equation}
\label{eq:Psi1-main-bound}
\big|
\Psi_1(h)-(1-\gamma h)
\big|
\leq
\big(
\tfrac{\gamma^2}{2}
+
\kappa
\big)h^2.
\end{equation}
Similarly, since $|\psi(r)-1|\leq\gamma r$ for $r\geq0$, Assumption \ref{as:coeff-approx-general} implies that, for $0\leq r\leq h$,
\begin{equation}
\label{eq:Psi2-main-bound}
\big|
\Psi_2(r)-1
\big|
\leq
(\gamma+\kappa)h.
\end{equation}
Therefore, \eqref{eq:velocity-update} can be rewritten as
\begin{equation}
\label{eq:velocity-decomposition-new}
\begin{aligned}
\bar{V}_{k+1}
=&
(1-\gamma h)\bar{V}_k
-
\alpha h\nabla U(\bar{X}_k)
+
J_1(h)\bar{V}_k
-
\alpha h
J_2(h-\tau_kh)
\nabla U(\bar{X}_{k+\tau_k})
\\    &
-
\alpha h
\Big(
\nabla U(\bar{X}_{k+\tau_k})
-
\nabla U(\bar{X}_k)
\Big)
+
\sqrt{2\gamma\alpha}
J_3(h),
\end{aligned}
\end{equation}
where we denote
\begin{equation}
\label{eq:velocity-J-def}
\begin{aligned}
J_1(h)
:=&
\Psi_1(h)-(1-\gamma h),
\\
J_2(h-\tau_kh)
:=&
\Psi_2(h-\tau_kh)-1,
\\
J_3(h)
:=&
\Delta W_{k+1}^{\Psi_3,2}.
\end{aligned}
\end{equation}
%%%
%%%
%%
By \eqref{eq:Psi1-main-bound} and \eqref{eq:Psi2-main-bound}, we have
\begin{equation}
\label{eq:velocity-I-basic-bound}
|J_1(h)|
\leq
\big(
\tfrac{\gamma^2}{2}
+
\kappa
\big)h^2,
\qquad
|J_2(h-\tau_kh)|
\leq
(\gamma+\kappa)h.
\end{equation}
Moreover, by It\^o's isometry and \eqref{eq:bound-Psi-Phi},
\begin{equation}
\label{eq:velocity-I3-bound}
\mathbb{E}
\big[
J_3(h)
\mid
\mathcal F_{t_k}
\big]
=0,
\qquad
\mathbb{E}
\Big[
\big|
J_3(h)
\big|^2
\Big]
\leq
(1+\kappa)^2dh.
\end{equation}
%
Taking square on both sides of \eqref{eq:velocity-decomposition-new}, taking expectations and using Young's inequality yield
\begin{equation}
\label{eq:velocity-square-start}
\begin{aligned}
&
\tfrac{1}{4}
\mathbb{E}
\Big[
\big|
\bar{V}_{k+1}
\big|^2
\Big]
\\ 
\leq&
\underbrace{
\tfrac{1}{4}
\mathbb{E}
\Big[
\big|
(1-\gamma h)\bar{V}_k
-
\alpha h\nabla U(\bar{X}_k)
\big|^2
\Big]
}_{=:\mathcal J_1}
+
\underbrace{
\gamma  \alpha
\mathbb{E}
\Big[
\big|
J_3(h)
\big|^2
\Big]
}_{=:\mathcal J_2}
+
\underbrace{
\tfrac{1}{4}h^2
\mathbb{E}
\Big[
\big|
(1-\gamma h)\bar{V}_k
-
\alpha h\nabla U(\bar{X}_k)
\big|^2
\Big]
}_{=:\mathcal J_3}
\\
&+
\underbrace{
\tfrac{1}{4}h^{-2}
\mathbb{E}
\Big[
\big|
J_1(h)\bar{V}_k
-
\alpha hJ_2(h-\tau_kh)
\nabla U(\bar{X}_{k+\tau_k})
-
\alpha h
\big(
\nabla U(\bar{X}_{k+\tau_k})
-
\nabla U(\bar{X}_k)
\big)
\big|^2
\Big]
}_{=:\mathcal J_4}
\\
&+
\underbrace{
\tfrac{1}{2}
\mathbb{E}
\Big[
\big|
J_1(h)\bar{V}_k
-
\alpha hJ_2(h-\tau_kh)
\nabla U(\bar{X}_{k+\tau_k})
-
\alpha h
\big(
\nabla U(\bar{X}_{k+\tau_k})
-
\nabla U(\bar{X}_k)
\big)
\big|^2
\Big]
}_{=:\mathcal J_5}
,
\end{aligned}
\end{equation}
where some cross terms containing $J_3(h)$ vanish after taking conditional expectation, since $J_3(h)$ has zero conditional mean given $\mathcal F_{t_k}$ and the other part in the inner product is $\mathcal F_{t_k}$-measurable.
%%
%%
%
We now estimate the terms $\mathcal J_i,i=1,\cdots,5$, on the right-hand side of \eqref{eq:velocity-square-start}. 
%First, expanding the leading part gives
Noting
\begin{equation}
\label{eq:velocity-leading-expand}
\begin{aligned}
&
\big|
(1-\gamma h)\bar{V}_k
-
\alpha h\nabla U(\bar{X}_k)
\big|^2
\\
=&
(1-\gamma h)^2
|\bar{V}_k|^2
-
2\alpha h(1-\gamma h)
\big\langle
\bar{V}_k,
\nabla U(\bar{X}_k)
\big\rangle
+
\alpha^2h^2
|\nabla U(\bar{X}_k)|^2
\end{aligned}
\end{equation}
and $(1-\gamma h)^2=1-2\gamma h+\gamma^2h^2$, we employ Young's inequality to infer
\begin{equation}
\label{eq:velocity-leading-final}
\begin{aligned}
\mathcal J_1
\leq&
\tfrac{1}{4}
\mathbb{E}
\Big[
|\bar{V}_k|^2
\Big]
-
\tfrac{\gamma}{2}h
\mathbb{E}
\Big[
|\bar{V}_k|^2
\Big]
-
\tfrac{\alpha}{2}h
\mathbb{E}
\Big[
\big\langle
\bar{V}_k,
\nabla U(\bar{X}_k)
\big\rangle
\Big]
\\
&+
\big(
\tfrac{\gamma^2}{4}
+
\tfrac{\gamma\alpha}{4}
\big)h^2
\mathbb{E}
\Big[
|\bar{V}_k|^2
\Big]
+
\big(
\tfrac{\alpha^2}{4}
+
\tfrac{\gamma\alpha}{4}
\big)h^2
\mathbb{E}
\Big[
|\nabla U(\bar{X}_k)|^2
\Big]
\\
\leq&
\tfrac{1}{4}
\mathbb{E}
\Big[
|\bar{V}_k|^2
\Big]
-
\tfrac{\gamma}{2}h
\mathbb{E}
\Big[
|\bar{V}_k|^2
\Big]
-
\tfrac{\alpha}{2}h
\mathbb{E}
\Big[
\big\langle
\bar{V}_k,
\nabla U(\bar{X}_k)
\big\rangle
\Big]
\\
&+
\big(
\tfrac{\gamma^2}{4}
+
\tfrac{\gamma\alpha}{4}
\big)h^2
\mathbb{E}
\Big[
|\bar{V}_k|^2
\Big]
+
\tfrac{\alpha^2+\gamma\alpha}{2}L^2h^2
\mathbb{E}
\Big[
|\bar{X}_k|^2
\Big]
+
\tfrac{\alpha^2+\alpha\gamma}{2}L_2^2dh^2,
\end{aligned}
\end{equation}
where \eqref{eq:gradient-growth-square} was used for the last step.
%%
%%
For the second term $\mathcal J_2$ in \eqref{eq:velocity-square-start}, we rely on \eqref{eq:velocity-I3-bound} to get
\begin{equation}
\label{eq:velocity-noise-final}
\mathcal J_2
\leq
\gamma\alpha(1+\kappa)^2dh.
\end{equation}
%%
%%
In view of \eqref{eq:gradient-growth-square} and the assumption $h\leq1$, we deduce
\begin{equation}
\label{eq:velocity-h2-leading-estimate}
\begin{aligned}
\mathcal J_3
\leq
\tfrac{1}{2}h^2
\mathbb{E}
\Big[
|\bar{V}_k|^2
\Big]
+
\alpha^2L^2h^2
\mathbb{E}
\Big[
|\bar{X}_k|^2
\Big]
+
\alpha^2L_2^2dh.
\end{aligned}
\end{equation}
%%
%%
It remains to estimate $\mathcal J_4, \mathcal J_5$ in \eqref{eq:velocity-square-start}. Invoking \eqref{eq:elementary_inequality}, 
\eqref{eq:grad-Xtau-split-clear}, \eqref{eq:velocity-I-basic-bound} and Lemma \ref{lem:nabla-tau-delta} yields
\begin{equation}
\label{eq:velocity-remainder-estimate}
\begin{aligned}
&\mathbb{E}
\Big[
\big|
J_1(h)\bar{V}_k
-
\alpha hJ_2(h-\tau_kh)
\nabla U(\bar{X}_{k+\tau_k})
-
\alpha h
\big(
\nabla U(\bar{X}_{k+\tau_k})
-
\nabla U(\bar{X}_k)
\big)
\big|^2
\Big]
\\
\leq&
3
\big(
\tfrac{\gamma^2}{2}
+
\kappa
\big)^2
h^4
\mathbb{E}
\Big[
|\bar{V}_k|^2
\Big]
+
3\alpha^2(\gamma+\kappa)^2h^4
\mathbb{E}
\Big[
|\nabla U(\bar{X}_{k+\tau_k})|^2
\Big]
\\  &
+
3\alpha^2h^2
\mathbb{E}
\Big[
\big|
\nabla U(\bar{X}_{k+\tau_k})
-
\nabla U(\bar{X}_k)
\big|^2
\Big]
\\
\leq&
\Big(
3\alpha^2(\gamma+\kappa)^2
\big(
4L^2+2\H_1^x
\big)
+
3\alpha^2\H_1^x
\Big)
h^4
\mathbb{E}
\Big[
|\bar{X}_k|^2
\Big]
\\
&+
\Big(
3
\big(
\tfrac{\gamma^2}{2}
+
\kappa
\big)^2
+
6\alpha^2(\gamma+\kappa)^2\H_1^v
+
3\alpha^2\H_1^v
\Big)
h^4
\mathbb{E}
\Big[
|\bar{V}_k|^2
\Big]
\\
&+
\Big(
3\alpha^2(\gamma+\kappa)^2
\big(
4L_2^2+2\H_1
\big)
+
3\alpha^2\H_1
\Big)dh^4,
\end{aligned}
\end{equation}
and therefore 
\begin{equation}
\label{eq:velocity-remainder-final}
\begin{aligned}
\mathcal J_4  + \mathcal J_5
\leq&
\Big(
\tfrac{9}{2}
\alpha^2(\gamma+\kappa)^2
\big(
2L^2+\H_1^x
\big)
+
\tfrac{9}{4}
\alpha^2\H_1^x
\Big)h^2
\mathbb{E}
\Big[
|\bar{X}_k|^2
\Big]
\\
&+
\Big(
\tfrac{9}{4}
\big(
\tfrac{\gamma^2}{2}
+
\kappa
\big)^2
+
\tfrac{9}{2}
\alpha^2(\gamma+\kappa)^2\H_1^v
+
\tfrac{9}{4}
\alpha^2\H_1^v
\Big)
h^2
\mathbb{E}
\Big[
|\bar{V}_k|^2
\Big]
\\
&+
\Big(
\tfrac{9}{2}
\alpha^2(\gamma+\kappa)^2
\big(
2L_2^2+\H_1
\big)
+
\tfrac{9}{4}
\alpha^2\H_1
\Big)
dh^2.
\end{aligned}
\end{equation}
Finally, substituting \eqref{eq:velocity-leading-final}-\eqref{eq:velocity-h2-leading-estimate} and \eqref{eq:velocity-remainder-final} into \eqref{eq:velocity-square-start} and noting $h\leq1$, we arrive at \eqref{eq:velocity-one-step}. This completes the proof.
\end{proof}

\end{document}
